\documentclass[11pt,a4paper]{article}

\usepackage[utf8]{inputenc}
\usepackage[margin=1in]{geometry}
\usepackage{amsmath,amssymb,bm}
\usepackage{booktabs}
\usepackage{hyperref}
\usepackage{xcolor}

\hypersetup{
  colorlinks=true,
  linkcolor=blue!60!black,
  urlcolor=blue!60!black,
  citecolor=blue!60!black
}

\newcommand{\state}{\bm{x}}
\newcommand{\meas}{\bm{z}}
\newcommand{\gain}{K}
\newcommand{\mask}{\bm{m}}
\newcommand{\avail}{\bm{a}}
\newcommand{\wrap}{\operatorname{wrap}}
\newcommand{\diag}{\operatorname{diag}}

\title{Small Report on Robust KalmaNet Trust for QCar}
\author{Quang Huy Nugyen}
\date{\today}

\begin{document}
\maketitle

\begin{abstract}
This report summarizes the Robust KalmaNet estimator implemented for the QCar
platform. The estimator keeps the physical QCar prediction model and learns a
trust mechanism for the measurement update. The trust mechanism is represented
by a learned measurement mask and a learned Kalman-style gain. Its purpose is to
reduce the influence of corrupted GPS position, heading, and velocity
measurements while preserving normal Extended Kalman Filter (EKF)-like behavior
on clean data.
\end{abstract}

\section{Motivation}

The QCar observer receives measurements from GPS, IMU, steering, wheel-speed,
and motor tachometer signals. A classical EKF can produce good estimates when
the measurement noise model is correct, but it can also over-trust a sensor
during abnormal events such as GPS jumps, GPS dropout, frozen velocity, wheel
speed scaling, or noisy inertial data.

The Robust KalmaNet design addresses this issue by separating the estimation
problem into two parts:
\begin{enumerate}
  \item a physics-based prediction step that propagates the vehicle state using
  the QCar kinematic model;
  \item a learned update step that decides how much to trust each measurement
  channel before applying a Kalman-style correction.
\end{enumerate}

The key idea is not to replace the vehicle model with a black-box neural
network. Instead, learning is concentrated on the part where uncertainty and
sensor faults are most difficult to model manually: measurement trust.

\section{State, Measurements, and Raw Inputs}

The active training surface uses the four-state vector
\begin{equation}
  \state_k =
  \begin{bmatrix}
    x_k & y_k & \psi_k & v_k
  \end{bmatrix}^{T},
\end{equation}
where $x$ and $y$ are global position, $\psi$ is heading, and $v$ is forward
velocity. Some recorded files contain an additional yaw-rate channel, but the
training data builder slices the target and measurement arrays to the first
four channels.

The measurement vector is
\begin{equation}
  \meas_k =
  \begin{bmatrix}
    x_{gps,k} & y_{gps,k} & \psi_{fused,k} & v_{tach,k}
  \end{bmatrix}^{T},
  \qquad
  \meas_k = H\state_k + \bm{n}_k,
  \qquad
  H \approx I_4 .
\end{equation}

The raw branch inputs used by the model are grouped as:
\begin{itemize}
  \item IMU branch: acceleration and yaw-rate channels;
  \item steering branch: steering angle and related motion context;
  \item wheel branch: four wheel-speed channels;
  \item availability metadata: GPS valid flag, heading valid flag, GPS hold
  flag, and GPS age.
\end{itemize}

\section{Prediction Model}

The current configuration uses \texttt{predictor\_mode: kinematic}, so the
prediction step is analytical and has no learned parameters. The vehicle is
propagated using a bicycle-model yaw-rate approximation:
\begin{equation}
  \dot{\psi}_k = \frac{v_k \tan(\delta_k)}{L},
\end{equation}
where $\delta_k$ is the steering angle and $L$ is the wheelbase. The prior state
is
\begin{align}
  x_k^- &= x_{k-1}^+ + v_k \cos(\psi_{k-1}^+) \Delta t, \\
  y_k^- &= y_{k-1}^+ + v_k \sin(\psi_{k-1}^+) \Delta t, \\
  \psi_k^- &= \wrap\left(
    \psi_{k-1}^+ + \frac{v_k \tan(\delta_k)}{L}\Delta t
  \right).
\end{align}

The active longitudinal option is \texttt{velocity\_lag}. In this mode,
velocity is predicted from throttle using
\begin{equation}
  \dot{v}_k = \frac{-v_{k-1}^+ + g u_k}{\tau},
  \qquad
  v_k^- = \operatorname{clip}\left(v_{k-1}^+ + \dot{v}_k\Delta t\right).
\end{equation}
The current config uses $L=0.2$, $\tau=0.301$, gain $g=6.598$, maximum velocity
$2.0$ m/s, and maximum acceleration $2.0$ m/s$^2$.

\section{Learned Trust Update}

The learned update receives the prior state $\state_k^-$, the previous updated
state $\state_{k-1}^+$, the current measurement $\meas_k$, the previous
measurement $\meas_{k-1}$, and sensor availability metadata.

The innovation and measurement difference are
\begin{align}
  \bm{r}_k &= \wrap(\meas_k - H\state_k^-), \\
  \Delta \meas_k &= \wrap(\meas_k - \meas_{k-1}).
\end{align}

The hard availability vector is
\begin{equation}
  \avail_k =
  \begin{bmatrix}
    g_k & g_k & h_k & 1
  \end{bmatrix}^{T},
\end{equation}
where $g_k$ is the GPS-position availability flag and $h_k$ is the heading
availability flag. The update feature vector combines normalized state
increment, innovation, measurement increment, hard availability, GPS hold flag,
GPS age, and innovation magnitudes. For the four-state model this produces a
26-dimensional update feature vector.

The trust mask is generated by a small MLP:
\begin{equation}
  \mask_k = \sigma(MLP(\bm{f}_k)).
\end{equation}
The measurement part of this vector, $\mask_k^z$, is interpreted as learned
measurement trust. Clean channels should have values near one. Attacked or
unreliable channels should move toward zero.

The Kalman-style gain is generated by a GRU and an MLP:
\begin{align}
  \bm{h}_k &= GRU(\mask_k \odot \bm{f}_k, \bm{h}_{k-1}), \\
  \gain_k &= s\tanh(MLP(\bm{h}_k)).
\end{align}
The current tanh scale is $s=1.2$. Both the measurement mask and gain are
smoothed with exponential moving averages using $\alpha=0.35$.

The final robust correction is
\begin{equation}
  \state_k^+
  =
  \state_k^-
  +
  \gain_k
  \left(
    \avail_k \odot \mask_k^z \odot \bm{r}_k
  \right).
\end{equation}
This equation is the main trust mechanism. A large raw innovation from a bad GPS
or velocity measurement cannot directly force a large state jump if the learned
mask suppresses that channel. The implementation also limits per-step update
corrections with the configured clamp
\[
  [0.35,\ 0.35,\ 0.35,\ 0.80]
\]
for $[x,y,\psi,v]$.

\section{Dataset and Training Procedure}

The dataset recorder saves synchronized samples into compressed \texttt{.npz}
files and matching JSON metadata. The active training configuration points to
\path{robust_kalmannet_dataset_V0_20260417_094048.npz}. Its metadata
reports 52,921 samples, 671.84 seconds of driving, and a mean timestep of
0.0127 seconds.

The trainer uses sliding windows with length $T=50$ for updater training and
$T=80$ in Phase C fine tuning. The validation split is the final 20 percent of
the sequence. Measurements are rebuilt offline using the local QCar EKF heading
path, and invalid GPS $x/y$ samples use the configured \texttt{freeze} behavior.

Because the predictor is kinematic, Phase A is automatically skipped. Training
therefore proceeds as:
\begin{itemize}
  \item Phase B: freeze the predictor and train the learned updater;
  \item Phase C: fine tune the full runtime loop with teacher forcing disabled.
\end{itemize}

The state loss uses weighted MSE with weights
\[
  W = [1.0,\ 1.0,\ 2.5,\ 1.2].
\]
The combined objective is
\begin{equation}
\begin{aligned}
  \mathcal{L}
  &=
  \lambda_u L_{upd}
  + \lambda_p L_{pred}
  + \lambda_m L_{mask}
  + \lambda_K L_K
  + \lambda_{\Delta K}L_{\Delta K}
  + \lambda_{\Delta m}L_{\Delta m}.
\end{aligned}
\end{equation}
Here $L_{upd}$ trains the final corrected state, $L_{pred}$ keeps the prior
reasonable, $L_{mask}$ supervises measurement trust under synthetic attacks,
$L_K$ penalizes large gains, and the temporal terms reduce gain and mask
chatter. Current Phase C weights are $\lambda_u=0.85$ and $\lambda_p=0.15$.

\section{Sensor Attack Augmentation}

Robustness is trained by corrupting raw branches and measurement channels. The
active augmentation mode is \texttt{mixed}. It includes:
\begin{itemize}
  \item raw branch attacks on IMU, steering, and wheel-speed branches;
  \item direct measurement attacks on heading and velocity;
  \item GPS position attacks on $x/y$.
\end{itemize}

The enabled raw attack types are bias, scale, freeze, noise, ramp, and zero-out.
GPS-specific attacks include noise, freeze, jump, dropout, and reacquisition
offsets. The active probabilities are:
\begin{center}
\begin{tabular}{lc}
\toprule
Attack group & Probability \\
\midrule
Raw branch attack & 0.25 \\
Direct measurement attack & 0.35 \\
GPS position attack & 0.45 \\
\bottomrule
\end{tabular}
\end{center}
Only one raw branch is attacked at a time. The supervision target for the
measurement trust mask is one for clean measurement channels and zero for
attacked channels.

\section{Observed Training Result}

The saved training history in \texttt{models/robust\_kalmannet.train\_history.json}
contains 20 logged epochs after Phase A was skipped. The best robust checkpoint
was selected at epoch 15 in Phase B.

\begin{center}
\begin{tabular}{lccc}
\toprule
Checkpoint point & Clean validation & Attacked validation & Selection loss \\
\midrule
Epoch 1 & 0.00668 & 1.10951 & 0.66838 \\
Best robust, epoch 15 & 0.01751 & 0.40279 & 0.24868 \\
Final Phase C, epoch 20 & 0.02408 & 0.43457 & 0.27037 \\
\bottomrule
\end{tabular}
\end{center}

The clean validation loss increases after aggressive attack training, but the
attacked validation loss decreases strongly. This is the expected tradeoff for a
robust-trust estimator: the model becomes more conservative on measurements so
that it can reject corrupted sensor channels.

\section{Conclusion}

The Robust KalmaNet trust mechanism for QCar combines a deterministic kinematic
prediction model with a learned measurement trust mask and a learned
Kalman-style gain. The predictor preserves vehicle physics, while the updater
learns when GPS, heading, and velocity should be trusted. The saved training run
shows that attack-aware training substantially reduces attacked validation loss,
which supports the intended use case: stable QCar state estimation under sensor
faults and GPS disturbances.

\end{document}
